\subsection*{Parte 1: Apuntar en una dirección (la tabla)}

\begin{center}
\renewcommand{\arraystretch}{1.3}
\begin{tabular}{P{2.8cm} P{6.0cm} P{3.0cm} P{2.8cm}}
\toprule
\textbf{Dirección en Scratch} & \textbf{Código de Python} & \textbf{Dirección en Python} & \textbf{¿Hacia dónde?} \\
\midrule
90 & \py|setheading(90 - 90)| & 0 & derecha \\
0 & \py|setheading(90 - 0)| & 90 & arriba \\
-90 & \py|setheading(90 - -90)| & 180 & izquierda \\
45 & \py|setheading(90 - 45)| & 45 & arriba y a la derecha \\
135 & \py|setheading(90 - 135)| & $-45$ & abajo y a la derecha \\
\bottomrule
\end{tabular}
\end{center}

Restar un número negativo es lo mismo que sumar, así que $90 - (-90) = 180$. Python también entiende $-45$: es la misma dirección que 315.

\subsection*{Parte 1: Tradúcelo tú}

\begin{code}
from ColabTurtlePlus.Turtle import *
clearscreen()
initializeTurtle((480, 360))
penup()
goto(-100, 0)
pendown()
pencolor("red")
pensize(4)
setheading(90 - 0)
forward(80)
right(90)
forward(80)
stamp()
\end{code}

La tortuga dibuja una línea roja hacia arriba desde $(-100, 0)$, gira a la derecha, dibuja una línea de lado a lado y se sella en $(-20, 80)$. Escribir \py|setheading(90)| en lugar de \py|setheading(90 - 0)| también funciona. Si agregaste \py|shape("turtle")|, ¡muy bien!: tu sello se ve como una tortuga.

\subsection*{Parte 1: Acertijo de dos líneas}

\begin{code}
setheading(90 - 180)
forward(100)
\end{code}

\py|setheading(-90)| también funciona, y también \py|right(90)| en lugar de \py|setheading|, porque la tortuga empieza mirando a la derecha. Un cuarto de vuelta en el sentido de las agujas del reloj la deja mirando hacia abajo.
